\frontmatterentry{致读者}

\begin{enumerate}
\def\labelenumi{\arabic{enumi}.}
\item
  《数学原本》丛书从数学的开端讲起，并给出完整的证明。原则上，它不要求读者具备任何专门的数学知识，而只要求对数学推理有一定的熟悉以及一定的抽象思维能力。不过，它尤其面向那些至少已经较好地掌握了大学数学课程前一两年内容的学生。
\item
  我们采用的叙述方法是公理化的，并且通常从一般过渡到特殊。证明的需要使题材具有严格固定的顺序。因此，除非读者已经具备相当广泛的数学知识，否则某些考虑的用处不会立刻显现出来。
\item
  丛书分为若干部，每部又分为若干章。法文版中已全部或部分出版的各书列于下面。在有英译本时，括号中给出相应的英文书名。在本卷中，凡有英译版时引用均指英译版，否则指法文版。
\end{enumerate}

Théorie des Ensembles（集合论）记作 E (S)

Algèbre（代数 \(^{(1)}\)）--- A (A)

Topologie Générale（一般拓扑）--- TG (GT)

Fonctions d'une Variable Réelle --- FVR

Espaces Vectoriels Topologiques（拓扑向量空间）--- EVT (TVS)

Intégration --- INT

Algèbre Commutative（交换代数 \(^{(2)}\)）--- AC (CA)

Variétés Différentielles et Analytiques --- VAR

Groupes et Algèbres de Lie（李群与李代数 \(^{(3)}\)）--- LIE (LIE) Théories Spectrales --- TS

\pandocbounded{\includegraphics[keepaspectratio]{images/part001_af8bcccf.jpg}}

在前六书（按上述顺序）中，正文里的每个命题都只假定已知那些已经在同一章或前面各章中讨论过的结果，这些章的顺序如下：S；A，第 I 至 III 章；GT，第 I 至 III 章；A，第 IV 章及以后；GT，第 IV 章及以后；FVR；TVS；INT。

从第七书起，读者通常可以找到关于该书与其他各书之间逻辑关系的精确说明（前六书总是被视为已知）。

\begin{enumerate}
\def\labelenumi{\arabic{enumi}.}
\setcounter{enumi}{3}
\item
  不过，我们有时在正文中插入了这样一些例子，它们涉及读者可能已经知道、但在本丛书中尚未讨论的事实。这样的例子置于两个星号之间：*\ldots*。大多数读者无疑会发现这些例子有助于理解正文。在另一些情形，*\ldots* 之间的段落涉及丛书中别处讨论的结果。我们希望读者能够自行验证其中不存在任何恶性循环。
\item
  每章的逻辑骨架由该章的定义、公理和定理组成。这些是以后使用时主要需要记住的部分。较次要的结果以及那些容易由定理推出的结果，标为“命题”、“引理”、“推论”、“注记”等。初次阅读时可以略去的部分用小字排印。偶尔会以“评注”（scholium）的名义给出关于某个特别重要定理的评述。
\end{enumerate}

为避免令人生厌的重复，有时方便的做法是引入仅在某一章或某一章的某一节内有效的记号或缩写（例如，在只讨论交换环的一章中，“环”一词总是指“交换环”）。这类约定总会被明确指出，通常出现在采用它们的章的开头。

\begin{enumerate}
\def\labelenumi{\arabic{enumi}.}
\setcounter{enumi}{5}
\item
  正文中有些段落旨在提醒读者警惕严重的错误。这些段落以记号 \(\sum\)（“危险弯道”）在页边标出。
\item
  习题的设置，既是为了让读者检验自己是否已消化正文，也是为了让读者注意到那些虽不宜写入正文、但仍有意义的结果。最难的习题带有记号 \(\P\) 。
\item
  一般说来，我们遵循通行的术语，除非有充分的理由偏离它。
\item
  我们特别力求始终使用严格正确的语言，同时不牺牲简洁性。凡属语言滥用之处，我们尽可能在正文中予以指明；没有这类语言滥用，任何数学文本都有流于迂腐、乃至无法阅读的危险。
\item
  由于正文原则上是对理论的教条式阐述，它一般不包含对文献的引用。文献的引用集中在历史注记中。每篇历史注记后面所附的参考文献，一般只包括那些在所讨论理论的发展过程中起过最重要作用的书籍和原始论文。它绝无求全之意。
\end{enumerate}

至于习题，我们一般不认为有必要指明其出处，因为它们取自许多不同的来源（原始论文、教科书、习题集）。

\begin{enumerate}
\def\labelenumi{\arabic{enumi}.}
\setcounter{enumi}{10}
\tightlist
\item
  在本书中，对定理、公理、定义 \ldots{} 等的引用，依次按如下方式给出：
\end{enumerate}

--- 引用法文版时，给出书（使用第 3 节所列的缩写）、章以及所在页码；

--- 引用本书时，只给出章和页码；

--- 引用英文版时，给出章、节和小节。

结果摘要以字母 R 引用；例如“集合论，R”表示“《集合论》的结果摘要”。

